% answers_short.tex —— 2007 年·普通高等学校招生考试·山东卷·文科 · 简短答案速查
% 由 tools/make_short_answers.py 自 answers.tex 的 \gksolhead/\gkanswer 提取，
% 并与 2007/2007_数学真题_山东文_答案要点.md 逐题比对后生成，勿手改（改源 MD 后重跑）。

\begin{gknotice}
\textbf{使用说明}：本节只给最终答案便于快速核对，完整推导见解析版。答案为编写者独立推导，\textbf{非官方标准答案}。
\end{gknotice}

\gkpart{选择题}
\begin{center}\begin{tabular}{ccccccccccccc}
\toprule
题号 & 1 & 2 & 3 & 4 & 5 & 6 & 7 & 8 & 9 & 10 & 11 & 12 \\
\midrule
答案 & B & B & D & A & C & B & C & A & B & A & B & D \\
\bottomrule
\end{tabular}\end{center}

\gkpart{填空题}
\begin{itemize}[leftmargin=2.6em]
  \item[\textbf{13.}] $\dfrac{1}{2007}$。
  \item[\textbf{14.}] $4$。
  \item[\textbf{15.}] $m\leqslant-5$。
  \item[\textbf{16.}] $(x-2)^2+(y-2)^2=2$。
\end{itemize}

\gkpart{解答题}
\begin{itemize}[leftmargin=2.6em]
  \item[\textbf{17.}] （1）$\dfrac18$；（2）$c=6$。
  \item[\textbf{18.}] （1）$a_n=2^{n-1}$；（2）$T_n=\dfrac32n(n+1)\ln2$。
  \item[\textbf{19.}] 甲电视台 $100$ 分钟、乙电视台 $200$ 分钟，最大收益 $70$ 万元。
  \item[\textbf{20.}] （1）证明见解析；（2）$E$ 为 $DC$ 的中点。
  \item[\textbf{21.}] $ab>0$ 时无极值点；$ab<0$ 时有唯一极值点，极值为 $-\dfrac b2+\dfrac b2\ln\left(-\dfrac{b}{2a}\right)$。
  \item[\textbf{22.}] （1）$\dfrac{x^2}{4}+\dfrac{y^2}{3}=1$；（2）定点 $\left(\dfrac27,0\right)$。
\end{itemize}
